% !TeX root = ../drafts/09-isa-programming.tex
\section{Programming Rust guests}\label{sec:isa-programming}

Guests are ordinary Rust \texttt{no\_std} executables compiled directly to RV64IM ELF. \path{crates/guest} provides the \texttt{leanvm\_guest} SDK, field/hash helpers, ECALL wrappers, entry macro, allocator, startup code, and linker script. There is no DSL compiler or instruction transpilation stage.

\subsection{Build and entry point}\label{sec:prog-functions}

The exact target is \texttt{riscv64im-unknown-none-elf}. It is a Tier 3 Rust target: use nightly Cargo with \texttt{rust-src} and build \texttt{core,alloc} from source. Installing a prebuilt target library is not a substitute. From the repository root:
\begin{verbatim}
rustup toolchain install nightly --component rust-src
./scripts/build-guests.sh
cargo build --release
\end{verbatim}
The script invokes the equivalent target build with \texttt{cargo +nightly}, \texttt{-Z build-std=core,alloc}, and \texttt{--target riscv64im-unknown-none-elf}, together with the guest linker/runtime flags. Host Cargo builds retain the native host target and can invoke this guest build through host integration. Do not set the workspace's default target to RISC-V. Targets containing A, C, floating-point, or CSR extensions do not describe this machine.

A guest executable enables the SDK's \texttt{runtime} feature, links with \path{crates/guest/link.x}, and defines one entry with \texttt{leanvm\_guest::entry!}. The entry function has type \texttt{fn() -> u64}; zero means successful EXIT, and a panic exits unsuccessfully. For example:
\begin{verbatim}
#![no_std]
#![no_main]
fn guest_main() -> u64 {
    let expected = leanvm_guest::public_input();
    let mut witness = [0u8; 32];
    leanvm_guest::read_witness_exact(&mut witness);
    if witness == expected { 0 } else { 1 }
}
leanvm_guest::entry!(guest_main);
\end{verbatim}
This example checks equality of witness bytes and public bytes; it does not by itself authenticate a signature or another application statement.

\paragraph{psABI and storage.} Calls use the RISC-V integer psABI: arguments in \texttt{a0..a7}, results in \texttt{a0/a1}, return address in \texttt{ra}, callee-saved registers as specified by the ABI, and a $16$-byte-aligned stack. The linker places the image at \texttt{0x10000}, separates executable and writable segments, and reserves a $16$ MiB stack below $2^{32}$. Startup initializes \texttt{gp}, \texttt{sp}, and BSS before calling the entry function. The heap ends below the stack reservation. Its bounded bump allocator does not reclaim storage on drop; allocations last until EXIT. Guest code must respect stack and heap limits.

\subsection{Untrusted witness data}\label{sec:prog-hints}

\texttt{read\_witness\_exact} reads from a sequential byte stream and fails on short input. Lengths, indices, cryptographic objects, and deferred child metadata read from it remain untrusted. Validate sizes before allocation and arithmetic, validate canonical encodings, and check every security-relevant claim before returning zero. Merely copying a witness to RAM is not verification. Mutable stores overwrite bytes; a second write does not assert equality with the first.

\subsection{Arithmetic and comparison}\label{sec:prog-div-ne}

Rust integer operations compile to native integer semantics. Use explicit wrapping operations when modular arithmetic is intended and checked operations where overflow must reject. Rust's source-level division checks can panic even though the underlying RV64IM DIV instruction has defined division-by-zero behavior. Equality, inequality, and integer range tests use ordinary Rust comparisons; no field inverse or exponent-address gadget is needed.

\subsection{Loops, branches, and matches}\label{sec:prog-loops}\label{sec:prog-conditionals}\label{sec:prog-match}

Loops maintain ordinary integer counters and mutable state. Branches and matches compile to RISC-V comparisons, branches, and compiler-selected control flow; function calls use the stack and psABI. Pointers are byte addresses. Bounds checks on witness-controlled indices must be preserved or replaced with an equivalent explicit check, not inferred from the fact that an address lies somewhere in RAM.

\subsection{Field values and range checks}\label{sec:prog-subfield-checks}\label{sec:prog-range-checks}

The SDK's \texttt{Field} represents an $\E$ value as three $64$-bit polynomial-basis limbs. Membership in $\K$ means the upper two limbs are zero, checked explicitly if the application needs it. Field addition is limbwise XOR; multiplication uses the proof-checked F192 ECALL on the guest target. A field element is not a machine pointer or an integer range proof. The CPU circuit already constrains native word/byte widths, while application-specific bounds remain guest checks.

\subsection{Cryptographic services}\label{sec:prog-byte-counter}

Use \texttt{Hasher} or \texttt{hash} for complete BLAKE2s hashing, and \texttt{blake2s\_compress} only when the application intentionally supplies its chaining value, counter, and flags. The wrapper calls the compression ECALL; the hash implementation handles message blocks, the integer byte counter, and finalization. \texttt{field\_mul} computes the three-limb extension-field product. Host builds use equivalent software primitives; guest builds issue the ABI of \S\ref{sec:rv-ecall}. These services prove primitive input/output relations, not application validity: the aggregation guest checks signatures, direct DA data, and coverage, while the field-native recursion circuit verifies proof transcripts.

\subsection{Recursive verifier integration}

\path{crates/riscv_proof} is \texttt{no\_std} with allocation support, but aggregation does not execute that verifier inside its guest. The shared SDK module \path{leanvm_guest::deferred} defines \path{Claim}, \path{public_fields}, \path{public_digest}, \path{claims_digest}, and \path{binding}. A \path{Claim} contains a statement digest, authorized key digest, and \texttt{u32} height. The aggregation guest uses \path{Statement::digest()} and \path{public_input(&statement, &children)} to bind the canonical statement and exact ordered child metadata to \path{leanvm_guest::public_input()}. Its \path{Child} contains the full canonical statement, key digest, and height, with neither ELF nor proof bytes. Successful \path{verify_input} establishes raw-claim validity and coverage conditional on those children; it does not authenticate their proofs.

On the host, construct \path{ProgramInfo::from_elf} from the trusted expected ELF and pass it to \path{recursion::Recursor::new}. Use \path{Recursor::key_digest()} when constructing the guest's child metadata, prove the RV64IM execution against its deferred binding, then call \path{Recursor::prove} with the statement digest, execution proof, children, and log inverse rate. Here the native \path{recursion::Child} instead contains the canonical child statement digest and a reference to its \path{NodeProof}. The returned \path{NodeProof} certifies the execution and all zero, one, or two children under the same authorized key. Root verification calls \path{Recursor::verify} with the statement digest and proof using the independently derived expected recursor. \path{NodeProof::to_bytes} and \path{NodeProof::from_bytes} provide strict versioned transport, not a replacement for verification (\S\ref{sec:recursion}).
